%%%%%     PACKS     %%%%%
\documentclass[12pt]{article}
\usepackage[margin=1in,headsep=.60in]{geometry}
\usepackage[utf8]{inputenc}
\usepackage[table, dvipsnames]{xcolor}
\usepackage{array}
\usepackage{amsmath}
\usepackage{booktabs}
\usepackage{mdframed} %For box around text
\usepackage{hyperref} %For hyperlinks in the PDF
\hypersetup{
    colorlinks=true,
    linkcolor=blue,
    filecolor=magenta,
    urlcolor=cyan,
}
\urlstyle{same}
\usepackage{amsthm}
\usepackage{amssymb}
\usepackage{amsfonts}
\usepackage{siunitx}
\usepackage{graphicx}
\usepackage{caption}
\usepackage{pgfplots}
\graphicspath{{Images/}}
\usepackage[colorinlistoftodos]{todonotes}
\usepackage{cleveref}
\usepackage[labelformat=simple]{subcaption}
\usepackage{grffile}
\usepackage{gensymb}
\usepackage{float}
\usepackage[shortlabels]{enumitem}
\usepackage{enumitem}
\setlistdepth{9}
\usepackage{biblatex} %Imports biblatex package
\addbibresource{references.bib}
\usepackage{outlines}
\usepackage{minted}
\usepackage{pdflscape}
\usepackage{everypage}
\usepackage{multirow}
\usepackage{multicol}
\usepackage{afterpage}
\usepackage[labelfont=bf, font=small]{caption}
\usepackage{upgreek}
\usepackage{tikz}
\usepackage{xcolor}
\usepackage{listings}

% style sheet for code
\definecolor{codegreen}{rgb}{0,0.6,0}
\definecolor{codegray}{rgb}{0.5,0.5,0.5}
\definecolor{codepurple}{rgb}{0.58,0,0.82}
\definecolor{backcolour}{rgb}{0.95,0.95,0.92}

\lstdefinestyle{mystyle}{
    backgroundcolor=\color{backcolour},   
    commentstyle=\color{codegreen},
    keywordstyle=\color{magenta},
    numberstyle=\tiny\color{codegray},
    stringstyle=\color{codepurple},
    basicstyle=\ttfamily\footnotesize,
    breakatwhitespace=false,         
    breaklines=true,                 
    captionpos=b,                    
    keepspaces=true,                 
    numbers=left,                    
    numbersep=5pt,                  
    showspaces=false,                
    showstringspaces=false,
    showtabs=false,                  
    tabsize=2
}

\lstset{style=mystyle}

\usepackage{lastpage}
%%%%%     COMMANDS     %%%%%
%%%% Blue box for subsection text (not figures) %%%%
\newenvironment{bluebox}
  {\begin{mdframed}[backgroundcolor=blue!5,linecolor=blue!40,roundcorner=8pt]}
  {\end{mdframed}}

%%%% Simple neutral box for figures %%%%
\newenvironment{figbox}
  {\begin{mdframed}[roundcorner=8pt,shadow=true,shadowsize=4pt,shadowcolor=black!40]}
  {\end{mdframed}}

\begin{document}

\begin{titlepage}
\newcommand{\HRule}{\rule{\linewidth}{0.5mm}}

\center

\textsc{\LARGE Aarhus University}\\[1.5cm]
\textsc{\Large Department Of Mathematics}\\[0.5cm]
\textsc{\large Introduction to Algebra 1}\\[0.5cm]
    

\HRule\\[0.4cm]
    \center 
    {\huge\bfseries  Handin 5}\\[0.4cm] % Title of your document
\HRule\\[1.5cm]

\begin{minipage}{0.4\textwidth}
        \begin{flushleft}
            \large
            \textit{Author}\\   
            Søren M. \textsc{Damsgaard}\\
                 % Your name
        \end{flushleft}
    \end{minipage}
~
    \begin{minipage}{0.4\textwidth}
        \begin{flushright}
            \large
            \textit{Student number}\\
            \textbf{202309814}\\
                 % Studienummer\
            
        \end{flushright}
    \end{minipage}

\vfill\vfill\vfill % Position the date 3/4 down the remaining page
    
    {\large\today}

\vfill\vfill
    \includegraphics[width=0.2\textwidth]{Aarhus_University_seal.png}\\[1cm] % Include a department/university logo - this will require the graphicx package
    

\vfill
\end{titlepage}
%%%%%      CHAPTERS      %%%%%
\hspace{0.02cm}
\section*{Goofy ah Math Number 5}
\section*{Digit sum!}
Let's remind ourselves what the digit sum is, so look at $12345$ that has the digit sum $1+2+3+4+5=15$ and generally the a number can be represented as
\begin{equation}\label{tsum}
    a_na_{n-1}\hdots a_1a_0 = \sum_{i=0}^{n}a_i10^i
\end{equation}
where the left-side represent the digits of the number, not some operation between some numbers $a_na_{n-1}\hdots a_1a_0$
\subsection*{Show some $n \in \mathbb{N}$ is $3|n$ iff. the digit sum of $n$, $n_{tsum}$ is also $3|n_{tsum}$}
This is equivalent to show
\[
n_{tsum} \equiv 0 \mod 3 \Leftrightarrow n \equiv 0 \mod 3
\]
start assuming that $n_{tsum} \equiv 0 \mod 3$ and rewriting $n_{tsum}$ using \eqref{tsum} and let
\begin{equation}
    n_{tsum} = a_n + a_{n-1} + \hdots + a_1+ a_0 
\end{equation}
the substitute $n_{tsum}$ with the above
\begin{equation}
    a_n + a_{n-1} + \hdots + a_1+ a_0 \equiv 0 \mod 3
\end{equation}
Then we want to show for all $a \in n_{Tsum}$ that $a\cdot 10^i \equiv a\cdot 10^k \mod 3$.\\
To do that let's look at properties of the number 10
\begin{equation*}
    10 \equiv 1 \mod 3
\end{equation*}
then remember that multiplication is preserved in congruence, so
\begin{equation*}
    10^i \equiv 1 \mod 3 \;\; \text{for } i\in\mathbb{N}
\end{equation*}
and then use lemma 9.2~\cite{Algebra} again, so
\begin{align*}
    a_i \cdot 10^i \equiv a_i \cdot 1 \mod 3\\
    =  a_i \cdot 10^i \equiv a_i \mod 3
\end{align*}
then apply this to every factor in $n_{Tsum}$ to get
\begin{align*}
    a_n\cdot 10^n &\equiv a_n \mod 3\\
    a_{n-1}\cdot10^{n-1} &\equiv a_{n-1} \mod 3\\
    &\vdots\\
    a_1 \cdot 10^1 &\equiv a_1 \mod 3\\
    a_0 \cdot 10^0 &\equiv a_0 \mod 3
\end{align*}
then use lemma 9.2~\cite{Algebra} again, to construct $n$
\begin{equation*}
    \sum_{i=0}^{n}a_i10^i \equiv \sum_{i=0}^{n}a_i\cdot 10^i \mod 3
\end{equation*}
which is
\begin{equation*}
    n \equiv a_n + a_{n-1} + \hdots + a_1+ a_0 \mod 3
\end{equation*}
\begin{equation*}
    n \equiv n_{Tsum} \mod 3
\end{equation*}
and remembering the original assumption $n_{Tsum} \equiv 0 \mod 3$ leads to the following
\begin{equation}
    n \equiv 0 \mod 3
\end{equation}
This can also be shown the other way around


\subsection*{Show some $n \in \mathbb{N}$ is $9|n$ iff. the digit sum of $n$, $n_{tsum}$ is also $9|n$}
The proof for this is equivalent to the above where you replace modulo 3 with modulo 9, since it holds that
\[
10 \equiv 1 \mod 9
\]
so using lemma 9.2~\cite{Algebra} results in the same conclusion.
\subsection*{let $n \in \mathbb{N}$ and let $m$ be a number that occures when you shift around the individual numbers in $n$. Show that $n-m$ is a multiple of $9$}
Let $n_{Tsum}$ be the digits sum of $n$ and let $m_{Tsum}$ be the digit sum of $m$. Since $m$ is formed by jumbling the digits of $n$ the digit sum of $m$ and $n$ is the same
\[
n_{Tsum} = m_{Tsum}
\]
and it's shown from the first part that a number $n$ is always congruent to its digit sum  modulo 9, this gives the following
\[
n \equiv n_{Tsum} \mod 9
\]
\[
m \equiv m_{Tsum} \mod 9
\]
and using the fact that their digit sums are the same
\[
n \equiv m \mod 9
\]
then using lemma 9.2.~\cite{Algebra} to subtract $m$ from both sides
\[
n-m \equiv 0 \mod 9
\]
which means the sum $n-m$ is a multiple of 9.
\printbibliography
\end{document}